\section{What existing sequence indexes store}
\label{sec:index-minimap2-comparison}

\subsection{Indexing: minimap2}

After computing minimizers they need to be grouped by their hash, minimap2 choose to use the lowest $b$ hash bits (by default $b=14$) to choose one of $2^b$ buckets~\autocite{minimap2_paper,minimap2_software}.
Minimap2 builds a hash table per bucket whose keys are the remaining hash bits. This buckets with hasmaps are serialized into \path{.mmi} file.
I performed in-depth analysis of how the \path{.mmi} format is specified based on available Github old issues~\autocite{minimap2_issue_820,minimap2_mmi_gist} and original code from minimap2's \texttt{mm\_idx\_dump} and \texttt{mm\_idx\_load} functions.
\noindent
\coderef{minimap2index/index.c}{Writing and loading an index part}{mm_idx_dump, mm_idx_load}{475}

The index file is concatenation of blocks presented in the Figure~\ref{fig:index-mmi-layout}:
\begin{itemize}
  \item \textbf{Magic bytes} The bytes are \texttt{MMI} followed by
    \texttt{0x02} (\texttt{MM\_IDX\_MAGIC}).
  \item \textbf{Metadata} Minimap2 saves $w$, the $k$-mer
    length $k$, used bitcount $b$, the number of sequences
    $n_{\mathrm{seq}}$ and a flag word (for various minimap2-specific bitflags)
  \item \textbf{Sequence table} For each of the $n_{\mathrm{seq}}$ sequences
    the file stores 1-byte name length $\ell$ followed by $\ell$ ASCII characters representing the sequence name (header taken from original FASTA file), and
    a 32-bit sequence length
  \item \textbf{Buckets} The buckets store the hash-map, where key is remaining hash and bit indicating if there's single value for that key or not. 
    The hash-map is stored as two components however. 
    First is the value list $p$ which is used when there are more than 1 value per key.
    Second is the hash array itself, which stores key/value pairs.
    If there are more than one value per key (last bit of key set to 1), then the value represents offset and count into the value list $p$
    The values, either as values in the key/value pairs or entries in $p$ store sequence id, position and strand bit.
  \item \textbf{Bit-acked bases (optional)} If the \texttt{MM\_I\_NO\_SEQ} flag
    is unset manually by user, bit-packed reference bases are saved at the end.
\label{lst:high-level-mmi-index-format-description}
\end{itemize}

The depiction of the index file layout is presented in the Figure~\ref{fig:index-mmi-layout} below.

\begin{figure}[htbp]
\centering
\begin{tikzpicture}[
  x=1cm,y=1cm,font=\footnotesize,
  bitfill/.style={draw=black!60,line width=0.4pt},
  lab/.style={align=center,font=\scriptsize,inner sep=0pt},
  zoom/.style={line width=0.9pt, line join=miter, line cap=rect,
    -{Latex[length=2.4mm,width=1.7mm]}}
]
  % (b) indent 1.05, (c)=(d) indent 2.10.
  % (b) ends at 8.55 so the green rail can pass at x=9.35.

  % ----- (a) y: boxes 14.20--14.95 -----
  \node[anchor=south west,font=\bfseries\footnotesize] at (0,15.25)
    {(a) Byte order of one \texttt{.mmi} part};
  \fill[indexBlue!14] (0,14.20) rectangle (1.70,14.95);
  \draw[black!55] (0,14.20) rectangle (1.70,14.95);
  \node[lab] at (0.85,14.575) {magic\\4\,B};
  \fill[indexBlue!14] (1.70,14.20) rectangle (5.20,14.95);
  \draw[black!55] (1.70,14.20) rectangle (5.20,14.95);
  \node[lab] at (3.45,14.575) {$w$, $k$, $b$, $n_{\mathrm{seq}}$, flag\\$5\times 4$\,B};
  \fill[indexOrange!16] (5.20,14.20) rectangle (7.80,14.95);
  \draw[black!55] (5.20,14.20) rectangle (7.80,14.95);
  \node[lab] at (6.50,14.575) {sequence table\\$n_{\mathrm{seq}}$ records};
  \fill[indexGreen!16] (7.80,14.20) rectangle (10.90,14.95);
  \draw[black!55] (7.80,14.20) rectangle (10.90,14.95);
  \node[lab] at (9.35,14.575) {buckets\\$2^{b}$ in order};
  \fill[black!8] (10.90,14.20) rectangle (12.80,14.95);
  \draw[black!55] (10.90,14.20) rectangle (12.80,14.95);
  \node[lab] at (11.85,14.575) {packed $S$\\if stored};

  % ----- (b) boxes 11.45--12.20, title below orange rail -----
  \node[anchor=south west,font=\bfseries\footnotesize] at (1.05,12.50)
    {(b) One sequence record};
  \fill[indexOrange!16] (1.05,11.45) rectangle (2.95,12.20);
  \draw[black!55] (1.05,11.45) rectangle (2.95,12.20);
  \node[lab] at (2.00,11.825) {name length $\ell$\\1\,B};
  \fill[indexOrange!16] (2.95,11.45) rectangle (6.45,12.20);
  \draw[black!55] (2.95,11.45) rectangle (6.45,12.20);
  \node[lab] at (4.70,11.825) {name (ASCII)\\$\ell$\,B};
  \fill[indexOrange!16] (6.45,11.45) rectangle (8.55,12.20);
  \draw[black!55] (6.45,11.45) rectangle (8.55,12.20);
  \node[lab] at (7.50,11.825) {seq.\ length\\4\,B};

  % ----- (c) boxes 8.75--9.50, title below green rail -----
  \node[anchor=south west,font=\bfseries\footnotesize] at (2.10,9.80)
    {(c) One bucket};
  \fill[indexGreen!16] (2.10,8.75) rectangle (3.70,9.50);
  \draw[black!55] (2.10,8.75) rectangle (3.70,9.50);
  \node[lab] at (2.90,9.125) {$n$\\4\,B};
  \fill[indexOrange!16] (3.70,8.75) rectangle (7.10,9.50);
  \draw[black!55] (3.70,8.75) rectangle (7.10,9.50);
  \node[lab] at (5.40,9.125) {value list $p$\\$n\times 8$\,B};
  \fill[indexGreen!16] (7.10,8.75) rectangle (8.60,9.50);
  \draw[black!55] (7.10,8.75) rectangle (8.60,9.50);
  \node[lab] at (7.85,9.125) {$s$\\4\,B};
  \fill[indexBlue!14] (8.60,8.75) rectangle (12.80,9.50);
  \draw[black!55] (8.60,8.75) rectangle (12.80,9.50);
  \node[lab] at (10.70,9.125) {$s$ key/value pairs\\$s\times(8{+}8)$\,B};

  % ----- (d) -----
  \node[anchor=south west,font=\bfseries\footnotesize] at (2.10,7.20)
    {(d) 64-bit words, high bits on the left ($b{=}14$)};

  \node[anchor=south west,font=\scriptsize] at (2.10,6.80) {Hash-table key};
  \fill[black!10] (2.10,6.05) rectangle (5.6109,6.65);
  \draw[bitfill] (2.10,6.05) rectangle (5.6109,6.65);
  \node[lab] at (3.86,6.35) {unused, 21 bits};
  \fill[indexBlue!18] (5.6109,6.05) rectangle (12.6328,6.65);
  \draw[indexBlue] (5.6109,6.05) rectangle (12.6328,6.65);
  \node[lab] at (9.12,6.35) {remaining hash $(\mathit{hash}\gg b)$, 42 bits};
  \fill[indexGreen!35] (12.6328,6.05) rectangle (12.80,6.65);
  \draw[indexGreen,very thick] (12.6328,6.05) rectangle (12.80,6.65);

  \node[anchor=south west,font=\scriptsize] at (2.10,5.50)
    {Hash-map value for single value / $p$ element};
  \fill[indexBlue!18] (2.10,4.70) rectangle (7.45,5.30);
  \draw[indexBlue] (2.10,4.70) rectangle (7.45,5.30);
  \node[lab] at (4.78,5.00) {sequence id, 32 bits};
  \fill[indexOrange!20] (7.45,4.70) rectangle (12.6328,5.30);
  \draw[indexOrange] (7.45,4.70) rectangle (12.6328,5.30);
  \node[lab] at (10.04,5.00) {last-base position, 31 bits};
  \fill[indexGreen!35] (12.6328,4.70) rectangle (12.80,5.30);
  \draw[indexGreen,very thick] (12.6328,4.70) rectangle (12.80,5.30);

  \node[anchor=south west,font=\scriptsize] at (2.10,4.20)
    {Hash-map value for non-single value};
  \fill[indexOrange!20] (2.10,3.40) rectangle (7.45,4.00);
  \draw[indexOrange] (2.10,3.40) rectangle (7.45,4.00);
  \node[lab] at (4.78,3.70) {offset into $p$, 32 bits};
  \fill[indexBlue!18] (7.45,3.40) rectangle (12.80,4.00);
  \draw[indexBlue] (7.45,3.40) rectangle (12.80,4.00);
  \node[lab] at (10.12,3.70) {count, 32 bits};

  \fill[indexGreen!35] (2.10,2.75) rectangle (2.38,3.03);
  \draw[indexGreen] (2.10,2.75) rectangle (2.38,3.03);
  \node[anchor=west,font=\scriptsize,text=black!55] at (2.52,2.89)
    {1-bit field: either means strand (for hash-map value) or means that hash-map value is single (for hash-map key)};

  % Orange (a seq -> b): H at mid(14.20, 12.20)=13.20. Spine 0.22.
  \draw[zoom,indexOrange]
    (6.50,14.20) -- (6.50,13.20) -- (0.22,13.20) -- (0.22,11.825) -- (1.05,11.825);
  % Green (a buckets -> c): H at mid(11.45, 9.50)=10.475. Spine 0.55.
  \draw[zoom,indexGreen]
    (9.35,14.20) -- (9.35,10.475) -- (0.55,10.475) -- (0.55,9.125) -- (2.10,9.125);
  % Blue (pairs -> key): H 7.95, below the p rail. Spine 1.25.
  \draw[zoom,indexBlue]
    (9.80,8.75) -- (9.80,7.95) -- (1.25,7.95) -- (1.25,6.35) -- (2.10,6.35);
  % Orange (p -> occurrence): H 8.40, above the blue rails. Spine 0.40.
  \draw[zoom,indexOrange]
    (5.40,8.75) -- (5.40,8.40) -- (0.40,8.40) -- (0.40,5.00) -- (2.10,5.00);
  % Blue (pairs -> value): right rail 13.60, enter from the right.
  \draw[zoom,indexBlue]
    (11.80,8.75) -- (11.80,7.95) -- (13.60,7.95) -- (13.60,3.70) -- (12.80,3.70);
  \draw[zoom,indexBlue]
    (11.80,8.75) -- (11.80,7.95) -- (13.60,7.95) -- (13.60,5.00) -- (12.80,5.00);
\end{tikzpicture}
\caption[Interpretation of one \texttt{.mmi} index]{%
Interpretation of one \texttt{.mmi} index.
For very large datasets index files are split into multiple parts where each of them follows the same format.
Minimap2 needs to store sequence ID (RID~\ref{lst:minimap-minimizer-cpp-decl}), because the index parts are generated for a set of sequences, not just one, so
each part needs to identify which reference sequence the minimizer comes from.
Arrows represent how specific elements of arrays look like in close inspection.
Panel~(a) depicts high-level structure of the index~\ref{lst:high-level-mmi-index-format-description}.
Panel~(b) presents records storing sequence information.
Panel~(c) shows each bucket of minimizers, which stores its own hash-map via two lists.
Panel~(d) shows how elements of $p$ looks like and key/values in the hash-map list.
Hash map key is always remaining 42 bits with bit indicating that there's single value for that key.
If the bit is set, then value in the key/value pair is (sequence id, position, strand bit), if not, then value points via offset/count pair into the list $p$.
The list $p$ stores the same (sequence id, position, strand bit) values.}
\label{fig:index-mmi-layout}
\end{figure}

The idea of storing hash-map as two lists inside bucket is also presented in the Figure~\ref{fig:index-mmi-construction} below to allow easier understanding of the 
index concept.

\begin{figure}[htbp]
\centering
\begin{tikzpicture}[x=1cm,y=1cm,font=\footnotesize]
  \node[anchor=west,font=\bfseries] at (0,5.8) {Minimizer occurrences before grouping};
  \foreach \x/\h/\p in {1/13/17,3/21/24,5/10/29,7/21/30}
    \node[index blue,minimum width=1.6cm] at (\x,4.95) {$h=\h$\\$p=\p$};
  \node[anchor=west,align=left] at (9,4.95) {$b=2$: bucket $h\bmod4$\\remaining key: $\lfloor h/4\rfloor$};
  \foreach \x/\b in {1/0,4/1,7/2,10/3}
    \node[index cell,fill=black!3,minimum width=2.2cm] (b\b) at (\x,3.3) {bucket \b};
  \draw[index arrow,indexBlue] (1,4.42)--(4,3.6);
  \draw[index arrow,indexBlue] (3,4.42)--(4,3.6);
  \draw[index arrow,indexBlue] (7,4.42)--(4,3.6);
  \draw[index arrow,black!45] (5,4.42)--(7,3.6);
  \node[anchor=west,font=\bfseries] at (0,2.35) {Bucket 1 after sorting};
  \node[index blue,minimum width=3.6cm] (single) at (2.35,1.5)
    {remaining bits 3: $p=17$};
  \node[index orange,minimum width=3.6cm] (repeated) at (2.35,.5)
    {remaining bits 5: offset 0, count 2};
  \node[anchor=west] at (6.0,1.7) {separate position array};
  \node[index orange,minimum width=1.15cm] at (6.65,.5) {$p=24$};
  \node[index orange,minimum width=1.15cm] at (7.85,.5) {$p=30$};
  \node[index cell,dashed,minimum width=1.15cm] at (9.05,.5) {$\cdots$};
  \node at (6.65,1.08) {0}; \node at (7.85,1.08) {1};
  \draw[index arrow,indexOrange] (repeated.east)--(6.05,.5);
  \node[anchor=west,text=black!65] at (0,-.45)
    {A singleton needs no separate list. Repeated keys share a contiguous array.};
\end{tikzpicture}
\caption{A simplified minimap2 hash bucket: hashes 13 and 21 have the same
lowest two bits, but different remaining keys, so they land in bucket 1 together. The two occurrences of hash
21 are placed next to each other in the bucket's position array.}
\label{fig:index-mmi-construction}
\end{figure}

\subsection{Indexing: mmseqs2}

There are alternative approaches in literature to building an index.
MMseqs2 stores one index file that contains counts and offsets, and then writes larger index files 
that contain all occurence lists instead of bucketing them like minimap2 does.
It's worth noting that 2017 MMseqs2\autocite{Steinegger079681,mmseqs2_2017} paper uses spaced k-mers which wasn't the focus of this thesis, so I didn't 
investigate more how the alignment is precisely done.

MMseqs2 offers great performance in clsutering, which is different problem to alignment, however connected as it involves 
estimation of edit distances for large datasets. The software also allows sets of taxonomy operations and alignment-based searches as well.\autocite{mmseqs2_linclust_2018}\autocite{mmseqs2_taxonomy_2021}\autocite{mmseqs2_desktop_2019}\autocite{mmseqs2_gpu_2025}.
\noindent
\coderef{mmseqsindex/src/prefiltering/IndexBuilder.cpp}{Count and fill the MMseqs2 index}{IndexBuilder::fillDatabase}{56}

\noindent
\coderef{mmseqsindex/src/prefiltering/IndexTable.h}{Read a word's entry range}{getDBSeqList}{178}

\noindent
\coderef{mmseqsindex/src/prefiltering/IndexTable.h}{Keep distinct word codes within one target}{addSequence}{350}

I illustrate the idea of MMseqs2 index on a small example in the Figure~\ref{fig:index-mmseqs-inverted} below, where 
I just assume that we use 2-letter k-mers of natural language for simplicity.

\begin{figure}[htbp]
\centering
\begin{tikzpicture}[x=1cm,y=1cm,font=\footnotesize]
  \node[anchor=west,font=\bfseries] at (0,6.55) {(a) Read target words};
  \foreach \x/\id/\s in {2/1/ACD,6/2/CAD,10/3/ACE}
    \node[index blue,minimum width=2.8cm] at (\x,5.75) {target \id: \texttt{\s}};
  \node[anchor=west,font=\bfseries] at (0,4.8) {(b) Count entries, then sum counts to find list starts};
  \node[anchor=east] at (1.25,3.95) {word};
  \node[anchor=east] at (1.25,3.2) {count};
  \node[anchor=east] at (1.25,2.45) {offset};
  \foreach \x/\word/\n/\o in {2.25/AC/2/0,4.1/AD/1/2,5.95/CA/1/3,7.8/CD/1/4,9.65/CE/1/5} {
    \node[index blue,minimum width=1.55cm] at (\x,3.95) {\texttt{\word}};
    \node[index cell,minimum width=1.55cm] at (\x,3.2) {$\n$};
    \node[index green,minimum width=1.55cm] at (\x,2.45) {$\o$};
  }
  \node[index green,minimum width=1.25cm] at (11.25,2.45) {6};
  \node[above] at (11.25,2.78) {end};
  \node[anchor=west,font=\bfseries] at (0,1.5) {(c) Fill the reserved list ranges};
  \foreach \x/\id/\p/\i in {2.1/1/0/0,3.7/3/0/1,5.3/2/1/2,6.9/2/0/3,8.5/1/1/4,10.1/3/1/5} {
    \node[index cell,minimum width=1.48cm] at (\x,.5) {$(\id,\p)$};
    \node[above] at (\x,.81) {$\i$};
  }
  \draw[indexBlue,very thick] (1.3,.12)--(4.5,.12);
  \node[below,text=indexBlue] at (2.9,.12) {AC: entries $[0,2)$};
  \node[anchor=west,align=left] at (6,-.15)
    {Entry = (target ID, word start).\\Positions start at zero.};
\end{tikzpicture}
\caption{A small example of MMseqs2's list construction, where word/count/offsets triplets are stored in a single file - practically only indices are stored there, while the actual (very large) lists of entries are stored separately.}
\label{fig:index-mmseqs-inverted}
\end{figure}
